% !TEX root = ../main.tex
\chapter{Phương pháp đề xuất}
\label{Chapter3}

Chương này trình bày chi tiết mô hình đề xuất, ký hiệu là SigLLM, cùng quy trình huấn
luyện và năm đóng góp của khóa luận. Xuất phát từ hai hạn chế đối lập được phân tích ở
Chương~\ref{Chapter2} -- mô hình ngôn ngữ lớn thuần văn bản thiếu tín hiệu lọc cộng tác,
trong khi mô hình lọc cộng tác thuần thiếu ngữ nghĩa -- SigLLM sử dụng một cầu nối
Q-Former để đưa biểu diễn lọc cộng tác vào mô hình ngôn ngữ lớn được giữ đông cứng.
Mục~\ref{sec:kien-truc-tong-quan} mô tả kiến trúc tổng thể. Các
Mục~\ref{sec:mf}--\ref{sec:ba-giai-doan} lần lượt trình bày ba khối thành phần và quy
trình huấn luyện ba giai đoạn. Các Mục~\ref{sec:cora}--\ref{sec:rank-loss} trình bày
năm đóng góp.

\section{Tổng quan kiến trúc}
\label{sec:kien-truc-tong-quan}

Bài toán được giải quyết là dự đoán tỉ lệ nhấp (Click-Through Rate -- CTR). Cho một
người dùng $u$, lịch sử tương tác của người dùng và một phim mục tiêu $i$, mô hình đưa
ra dự đoán nhị phân về khả năng người dùng thích phim mục tiêu, biểu diễn dưới dạng câu
trả lời ``Có'' hoặc ``Không''. Đây là bài toán phân loại nhị phân, được đánh giá bằng
hai độ đo AUC và uAUC (Chương~\ref{Chapter4}).

SigLLM gồm ba khối nối tiếp, minh họa ở Hình~\ref{fig:kien-truc}. Khối \textbf{Phân rã
ma trận} (Matrix Factorization -- MF) cung cấp các véc-tơ nhúng lọc cộng tác cho người
dùng và vật phẩm. Khối \textbf{cầu nối Q-Former} nén các véc-tơ này thành một khối
\textit{soft token} và căn chỉnh chúng với ngữ nghĩa văn bản của phim. Khối \textbf{mô
hình ngôn ngữ lớn} (Large Language Model -- LLM) thực hiện suy luận và sinh câu trả lời.

Một đặc điểm thiết kế của mô hình là phần lớn tham số được giữ đông cứng: cả MF lẫn LLM
đều không được cập nhật trong quá trình huấn luyện. Chỉ Q-Former, các lớp chiếu, bộ
chuyển thể LoRA~\cite{lora2021} và bộ tiêm CoRA được huấn luyện, với nhãn giám sát ở
giai đoạn cuối là một token ``Có/Không''. Cách thiết kế này bảo toàn tri thức ngôn ngữ
sẵn có của LLM và giảm chi phí huấn luyện.

\begin{figure}[htp]
\centering
\includegraphics[width=12cm]{images/logo.png}
\caption{Kiến trúc tổng quan của SigLLM. Khối MF đông cứng cung cấp tín hiệu lọc cộng
tác; Q-Former nén và căn chỉnh tín hiệu này thành soft token; LLM đông cứng thực hiện
suy luận. Chỉ Q-Former, các lớp chiếu, LoRA và bộ tiêm CoRA được huấn luyện.
(\textit{Hình minh họa tạm; cần thay bằng sơ đồ kiến trúc chính thức.})}
\label{fig:kien-truc}
\end{figure}

\section{Khối Phân rã ma trận}
\label{sec:mf}

Khối MF đóng vai trò nguồn tín hiệu lọc cộng tác. Từ ma trận tương tác người dùng--vật
phẩm, MF học ánh xạ mỗi người dùng $u$ và mỗi vật phẩm $i$ thành các véc-tơ nhúng ẩn
$e_u, e_i \in \mathbb{R}^{d}$ sao cho tích vô hướng $e_u^\top e_i$ phản ánh mức độ phù
hợp giữa người dùng và vật phẩm. MF được huấn luyện riêng theo mục tiêu xếp hạng cặp
BPR~\cite{bpr2009}, sau đó được đông cứng. Ba loại tín hiệu được trích ra để đưa vào cầu
nối: nhúng người dùng $e_u$, nhúng phim mục tiêu $e_i$ và nhúng các phim trong lịch sử
tương tác.

Việc huấn luyện MF tách rời khỏi các khối còn lại là một lựa chọn thiết kế nhằm cô lập
chất lượng của tín hiệu lọc cộng tác khỏi quá trình huấn luyện LLM, tạo điều kiện phân
tích rõ ràng vị trí nút thắt của hệ thống.

\section{Cầu nối Q-Former}
\label{sec:qformer}

Cầu nối được xây dựng dựa trên kiến trúc Q-Former theo InstructBLIP~\cite{instructblip2023},
vốn kế thừa từ BLIP-2~\cite{blip2_2023}. Q-Former sử dụng một tập nhỏ các truy vấn học
được. Trong mỗi lớp, các truy vấn đi qua một tầng self-attention cùng với các token của
chỉ dẫn (instruction) văn bản, sau đó cross-attention để thu nhận thông tin từ tín hiệu
lọc cộng tác do MF cung cấp. Đầu ra là một khối truy vấn $\mathbf{Z} = \text{cf\_q} \in
\mathbb{R}^{B \times Q \times d_{\text{model}}}$, với $Q$ là số truy vấn (mặc định
$Q=16$) và $d_{\text{model}}=768$.

Vai trò của các thành phần Q-Former thay đổi khi chuyển từ bối cảnh thị giác--ngôn ngữ
sang gợi ý. Trong InstructBLIP, cross-attention là cơ chế chính, thực hiện nén một chuỗi
dài các mảng ảnh thành một số ít truy vấn. Trong SigLLM, đầu vào của cross-attention chỉ
gồm một số ít véc-tơ lọc cộng tác, nên đường self-attention giữa truy vấn và chỉ dẫn trở
thành cơ chế trích xuất đặc trưng phụ thuộc tác vụ chủ yếu. Nhận định này giải thích quan
sát rằng việc tăng số lớp Q-Former không cải thiện kết quả, trong khi việc đa dạng hóa
chỉ dẫn lại có tác dụng.

Khối truy vấn $\mathbf{Z}$ được đưa vào LLM theo hai đường, trình bày chi tiết ở
Mục~\ref{sec:cora}: đường \textit{soft token} chiếu $\mathbf{Z}$ về không gian embedding
của LLM rồi chèn vào chỉ dẫn tại vị trí đánh dấu \verb|<CFTokens>|; và đường
\textit{trọng số} biến $\mathbf{Z}$ thành một lượng trọng số tăng thêm cộng vào các lớp
attention của LLM.

\section{Quy trình huấn luyện ba giai đoạn}
\label{sec:ba-giai-doan}

SigLLM được huấn luyện qua ba giai đoạn nối tiếp.

\textbf{Giai đoạn 1 --- Tiền huấn luyện biểu diễn.} Q-Former được huấn luyện bằng các mục
tiêu tương phản và căn chỉnh giữa tín hiệu lọc cộng tác và văn bản, bao gồm căn chỉnh
ảnh--văn bản (ITC), so khớp ảnh--văn bản (ITM), sinh văn bản có điều kiện (ITG), cùng các
mục tiêu tương phản vật phẩm--vật phẩm và người dùng--vật phẩm.

\textbf{Giai đoạn 2 --- Tiền huấn luyện sinh.} Q-Former và lớp chiếu được huấn luyện với
mục tiêu sinh, để đầu ra của cầu nối trở thành soft token mà LLM có thể diễn giải.

\textbf{Giai đoạn 3 --- Tinh chỉnh CTR.} Giai đoạn này theo giao thức hai bước của
CoLLM~\cite{collm2023}. Bước 1 huấn luyện LoRA cơ sở với chỉ dẫn chỉ gồm văn bản, không
có tín hiệu lọc cộng tác, thiết lập một đường cơ sở. Bước 2 huấn luyện mô hình đầy đủ với
tín hiệu lọc cộng tác được bơm vào (qua soft token, CoRA, hoặc cả hai), khởi tạo từ
checkpoint của Bước 1.

Năm đóng góp của khóa luận tác động vào các giai đoạn khác nhau và được tổng hợp ở
Bảng~\ref{tab:cai-tien}.

\begin{table}[ht]
\caption{Tổng hợp năm đóng góp của khóa luận và giai đoạn tác động}
\label{tab:cai-tien}
\begin{center}
\begin{tabular}{|c|L{5.5cm}|L{5.5cm}|}
\hline
\textbf{Mã} & \textbf{Đóng góp} & \textbf{Mục tiêu} \\
\hline
A & CoRA -- tiêm tín hiệu lọc cộng tác vào trọng số LLM & Đưa tín hiệu lọc cộng tác vào
sâu trong tính toán của LLM, tránh nhiễu ở không gian đầu vào \\
\hline
B & Tách lớp chiếu người dùng và vật phẩm & Ánh xạ đúng đặc trưng hình học của
user-factor và item-factor trong không gian MF \\
\hline
C & Backbone Instruct và chỉ dẫn ChatML & Tận dụng tiên nghiệm tuân theo chỉ dẫn cho bài
toán trả lời ``Có/Không'' \\
\hline
F & Cân bằng tương phản người dùng--vật phẩm ở Giai đoạn 1 & Bảo đảm mục tiêu căn chỉnh
vật phẩm--văn bản được huấn luyện đầy đủ \\
\hline
-- & Mục tiêu bảo toàn thứ tự & Cải thiện thứ tự xếp hạng nội-người-dùng (uAUC) \\
\hline
\end{tabular}
\end{center}
\end{table}

\section{Đóng góp A --- Tiêm tín hiệu lọc cộng tác vào trọng số LLM}
\label{sec:cora}

\textbf{Động cơ.} Cách tích hợp tín hiệu lọc cộng tác phổ biến, chẳng hạn
CoLLM~\cite{collm2023}, chèn soft token vào không gian đầu vào của chỉ dẫn. Khi đó tín
hiệu định danh và ngữ nghĩa văn bản cùng tồn tại trong một không gian, còn ảnh hưởng của
soft token tới các tầng sâu chỉ là gián tiếp. CoRA~\cite{cora2024} chỉ ra hai hạn chế của
hướng tiếp cận không gian đầu vào: tinh chỉnh LLM trên dữ liệu gợi ý làm suy giảm tri
thức ngôn ngữ, và việc đưa đặc trưng lọc cộng tác vào chỉ dẫn làm nhiễu ngữ nghĩa.

\textbf{Cơ chế.} Theo hướng của CoRA~\cite{cora2024}, khối truy vấn $\mathbf{Z} =
\text{cf\_q} \in \mathbb{R}^{B \times Q \times d_{\text{model}}}$ được biến thành một
lượng trọng số tăng thêm hạng thấp, riêng cho từng mẫu, cộng vào các lớp chiếu
$q_{\text{proj}}$ và $v_{\text{proj}}$ trong khối attention của LLM. Với đầu vào $x$ của
một lớp, lượng thay đổi $\Delta$ được tính theo Công thức~\ref{eq:cora-delta}:

\begin{equation} \label{eq:cora-delta}
\Delta = s \cdot g \cdot \big( (x \, A^\top) \, B \big),
\qquad s = \frac{\alpha}{Q},
\end{equation}

trong đó $A = P_{\text{down}}(\mathbf{Z})$ và $B = P_{\text{up}}(\mathbf{Z})$ là các phép
chiếu hạng thấp sinh từ $\mathbf{Z}$, $g$ là một cổng vô hướng, $s$ là hệ số tỉ lệ với
$\alpha$ là siêu tham số và $Q$ là số truy vấn. Đầu ra của lớp trở thành $\text{output}
+ \Delta$.

Điểm khác biệt so với LoRA tiêu chuẩn nằm ở nguồn của lượng thay đổi. LoRA học một lượng
thay đổi cố định, dùng chung cho mọi đầu vào, tức thích nghi theo tác vụ. Bộ tiêm ở đây
sinh $\Delta$ từ khối truy vấn $\mathbf{Z}$ của từng mẫu, nên mỗi người dùng và mỗi phim
tương ứng với một lượng thay đổi riêng. Bộ sinh $(P_{\text{down}}, P_{\text{up}})$ được
chia sẻ qua mọi tầng và chỉ tách theo dạng trọng số -- một bộ cho các lớp
$q_{\text{proj}}$ và một bộ cho các lớp $v_{\text{proj}}$ -- giúp giữ số tham số ở mức
thấp; toàn bộ biến thiên theo mẫu đến từ $\mathbf{Z}$.

\textbf{Bảo đảm ổn định.} Việc bổ sung một đường truyền mới lên một LLM đông cứng đòi hỏi
kiểm soát ổn định. Bốn kỹ thuật sau được áp dụng.

\begin{enumerate}
    \item \textit{Cổng tanh khởi tạo không}, theo Flamingo~\cite{flamingo2022}. Mỗi dạng
    trọng số có một cổng $g$ khởi tạo bằng $0$, do đó $\tanh(0)=0$ và $\Delta$ bằng $0$ tại
    thời điểm bắt đầu; LLM giữ nguyên hành vi ban đầu rồi điều chỉnh dần, với $\tanh$ giới
    hạn $\Delta$ trong khoảng $(-1,1)$. Thực nghiệm không có cổng này cho thấy mô hình phân
    kỳ và AUC suy giảm về mức ngẫu nhiên $0{,}5$.

    \item \textit{Không khởi tạo không cho $P_{\text{up}}$}. Nếu cả cổng lẫn
    $P_{\text{up}}$ đều bằng $0$ thì gradient của cả hai bằng $0$ và đường truyền không thể
    khởi động. Theo Flamingo, chỉ cổng được khởi tạo bằng $0$, còn phép biến đổi được khởi
    tạo với độ lệch chuẩn $0{,}02$.

    \item \textit{Chuẩn hóa truy vấn về chuẩn L2 đơn vị} trước bộ sinh, nhằm tách độ lớn
    của $\Delta$ khỏi chuẩn của truy vấn. Chuẩn của truy vấn quan sát được ở mức trung bình
    khoảng $16$, đủ lớn để gây bất ổn nếu không chuẩn hóa.

    \item \textit{Tính toán ở độ chính xác fp32 và xử lý đa thiết bị}, đồng thời bỏ qua
    khi kích thước lô không khớp trong quá trình giải mã tăng dần.
\end{enumerate}

\textbf{Cấu hình.} Bộ tiêm được điều khiển bởi chế độ tiêm (nhận một trong ba giá trị:
chỉ soft token, chỉ trọng số, hoặc cả hai), hệ số $\alpha$ và các mô-đun đích
$q_{\text{proj}}, v_{\text{proj}}$. Bộ tiêm được gắn sau LoRA, nên LoRA và CoRA cùng tồn
tại. Ngay cả ở chế độ chỉ trọng số, Q-Former vẫn hoạt động như bộ mã hóa lọc cộng tác,
vì $\mathbf{Z}$ là đầu ra của Q-Former; bộ tiêm chỉ xác định cách đọc đầu ra đó vào LLM.
Về chi phí, mỗi lần truyền xuôi bổ sung hai phép nhân ma trận trên mỗi lớp
$q_{\text{proj}}$ và $v_{\text{proj}}$, làm tăng thời gian và bộ nhớ; độ ổn định phụ thuộc
vào $\alpha$, nên giá trị này được tăng dần và theo dõi qua AUC trên tập kiểm định.

\section{Đóng góp B --- Tách lớp chiếu người dùng và vật phẩm}
\label{sec:proj}

Ở phiên bản cơ sở, một lớp chiếu chung ánh xạ cả nhúng người dùng lẫn nhúng vật phẩm từ
không gian MF sang không gian Q-Former. Tuy nhiên, đặc trưng hình học của véc-tơ
user-factor và véc-tơ item-factor trong không gian MF khác nhau, nên việc dùng chung một
phép chiếu là không phù hợp. Đóng góp B tách thành hai lớp chiếu riêng: một cho nhúng
người dùng và một cho các nhúng vật phẩm (phim mục tiêu và lịch sử). Nhờ đó mỗi loại
nhúng được ánh xạ theo đặc trưng riêng, cho biểu diễn lọc cộng tác đưa vào Q-Former
nhất quán hơn. Để duy trì tương thích, checkpoint có một lớp chiếu chung được nạp đồng
thời vào cả hai lớp mới.

\section{Đóng góp C --- Backbone Instruct và chỉ dẫn ChatML}
\label{sec:instruct}

Phiên bản cơ sở sử dụng backbone Qwen2-7B-Base, một mô hình không được tinh chỉnh theo
chỉ dẫn, với chỉ dẫn ở dạng hoàn thành văn bản. Đóng góp C thay bằng backbone
Qwen2-7B-Instruct~\cite{qwen2} và khôi phục định dạng ChatML tương ứng, gồm một thông
điệp hệ thống yêu cầu trả lời bằng ``Yes'' hoặc ``No'', một lượt người dùng chứa chỉ dẫn
cùng vị trí \verb|<CFTokens>|, và một lượt trợ lý. Bộ tách token được điều chỉnh để vị
trí soft token không trùng với các token điều khiển \verb+<|im_start|>+ và
\verb+<|im_end|>+.

Một phương án thay thế là tăng dung lượng LoRA trên backbone Base. Thực nghiệm cho thấy
phương án này làm giảm chất lượng: khi tăng hạng LoRA lên $r=16$ trên các lớp
$[q,k,v,o]$, mô hình quá khớp trên tập dữ liệu nhỏ, với mất mát huấn luyện giảm nhanh hơn
nhưng uAUC trên tập kiểm định giảm. Kết quả này cho thấy dung lượng LoRA không phải là
nút thắt; nút thắt nằm ở chất lượng căn chỉnh của Q-Former ở Giai đoạn 1. Điều còn thiếu
là một tiên nghiệm tuân theo chỉ dẫn tốt hơn, vốn được backbone Instruct cung cấp sẵn,
cho phép LoRA nhỏ tập trung vào việc sử dụng tín hiệu lọc cộng tác. Do backbone Instruct
được huấn luyện với định dạng ChatML còn backbone Base không xử lý các token điều khiển
tương ứng, việc thay backbone kéo theo thay đổi định dạng chỉ dẫn.

\section{Đóng góp F --- Cân bằng tương phản người dùng--vật phẩm ở Giai đoạn 1}
\label{sec:balance}

\textbf{Vấn đề.} Ở Giai đoạn 1, dữ liệu huấn luyện là một danh sách trộn ba loại mẫu:
vật phẩm--văn bản (mục tiêu căn chỉnh chính), vật phẩm--vật phẩm, và người dùng--vật
phẩm. Mỗi lô được tách theo loại và tính mất mát riêng, trong đó một mất mát chỉ được kích
hoạt khi lô chứa ít nhất hai mẫu cùng loại. Trên tập ML-1M, số cặp người dùng--vật phẩm
dương vào khoảng $888$ nghìn, trong khi số cặp vật phẩm--văn bản chỉ khoảng $3$ nghìn,
tức tỉ lệ mẫu vật phẩm--văn bản khoảng $0{,}34\%$. Với lô kích thước $64$, kỳ vọng số mẫu
vật phẩm--văn bản mỗi lô chỉ khoảng $0{,}22$, nên điều kiện ``ít nhất hai mẫu'' hiếm khi
được thỏa. Hệ quả là các mục tiêu căn chỉnh vật phẩm--văn bản (ITC, ITM, ITG) gần như
không được huấn luyện, làm suy giảm chất lượng căn chỉnh của Q-Former.

\textbf{Giải pháp.} Đóng góp F điều chỉnh ở hai tầng, tổng hợp ở Bảng~\ref{tab:balance}:
giới hạn số cặp người dùng--vật phẩm khi dựng dữ liệu, nhằm thay đổi tần suất xuất hiện;
và giảm trọng số mất mát người dùng--vật phẩm khi huấn luyện, nhằm điều hòa cường độ đóng
góp.

\begin{table}[ht]
\caption{Hai điều chỉnh của Đóng góp F}
\label{tab:balance}
\begin{center}
\begin{tabular}{|L{4cm}|L{3cm}|L{5cm}|}
\hline
\textbf{Điều chỉnh} & \textbf{Tầng} & \textbf{Tác động} \\
\hline
Giới hạn số cặp người dùng--vật phẩm ở mức $20{.}000$ & Dữ liệu & Nâng tỉ lệ mẫu
vật phẩm--văn bản lên khoảng $13\%$; với lô $64$, kỳ vọng khoảng $8$ mẫu, nên điều kiện
``ít nhất hai mẫu'' gần như luôn được thỏa \\
\hline
Giảm trọng số mất mát người dùng--vật phẩm từ $1{,}0$ xuống $0{,}5$ & Mất mát & Điều hòa
cường độ gradient người dùng--vật phẩm \\
\hline
\end{tabular}
\end{center}
\end{table}

Việc thay đổi giới hạn số cặp yêu cầu dựng lại tệp dữ liệu, vì giới hạn được cố định trong
tệp lúc dựng, trong khi trọng số mất mát được đọc lúc huấn luyện và có hiệu lực ngay. Giới
hạn hiện được thực hiện bằng cách lấy phần đầu của danh sách; nếu danh sách được sắp theo
người dùng hoặc thời gian thì cách này có thể gây thiên lệch, và một cải tiến là xáo trộn
danh sách trước khi cắt.

\section{Mục tiêu bảo toàn thứ tự cho xếp hạng nội-người-dùng}
\label{sec:rank-loss}

Bốn đóng góp trên chủ yếu cải thiện khả năng phân biệt ở mức tổng thể, tương ứng với độ
đo AUC. Độ đo uAUC đánh giá thứ tự xếp hạng trong phạm vi từng người dùng và bất biến với
các phép biến đổi đơn điệu riêng theo người dùng; do đó các tín hiệu chỉ khác biệt giữa
các người dùng không cải thiện uAUC. Cơ chế của hiện tượng này được phân tích ở
Mục~\ref{sec:nghich-ly} của Chương~\ref{Chapter4}. Nhằm cải thiện trực tiếp thứ tự
nội-người-dùng, khóa luận bổ sung một mục tiêu bảo toàn thứ tự (rank-preserving
collaborative-alignment) ở tầng căn chỉnh, thay vì bổ sung tín hiệu phân biệt giữa các
người dùng. Mục tiêu này là một thành phần tùy chọn và được đánh giá tách biệt.

\section{Kết chương}

Chương này đã trình bày kiến trúc ba khối của SigLLM, quy trình huấn luyện ba giai đoạn,
và năm đóng góp của khóa luận. Luận điểm xuyên suốt là nút thắt của hệ thống nằm ở chất
lượng căn chỉnh của Q-Former, chứ không phải dung lượng của mô hình ngôn ngữ lớn.
Chương~\ref{Chapter4} kiểm chứng các đóng góp này bằng thực nghiệm trên tập MovieLens-1M.
